\section{Executable evidence, reproducibility, and claim discipline}

\subsection{Reference implementation}

The archive contains \path{code/native_clause_framework.py}.  It implements,
with exact integer and rational arithmetic:

\begin{itemize}
\item bounded least search and complete enumeration;
\item least closure of a monotone transformer on a finite product of chains;
\item exhaustive monotonicity checks for small test boxes;
\item reverse demand propagation on a DAG;
\item clocked trace execution and invariant checking;
\item backward relational viability and least viable paths;
\item stable-observer diameter checks and motion tails;
\item exact label recovery from a separation gap;
\item conversion between a height bound and a finite list;
\item common good windows for finite exceptional sets;
\item lexicographic descent fuel;
\item execution of a finite contract graph with local host-level checkers.
\end{itemize}

The final item is deliberately documented as an executable analogue, not an
object proof.  The distinction is enforced in the module docstring.

\subsection{Regression tests}

The accompanying test suite has fourteen tests.  It covers success and failure
cases, exact rational stability, cyclic demand closure, demand minimality in a
small graph, trace invariants, relational failure, gap recovery, and graph
certificate rejection.  The recorded run in
\path{validation/python_tests.txt} reports all tests passing.

These tests provide three kinds of evidence:

\begin{enumerate}
\item the finite algorithms are specified precisely enough to execute;
\item examples in the report match the implementation;
\item later refactoring can detect changes in behavior.
\end{enumerate}

They do not prove the general theorems and do not check the Lean outline.

\subsection{Machine-readable schema}

\path{schema/native_clause_program.schema.json} describes a declarative program
format.  A node records its identifier, kind, dependencies, output formula
name, demand/resource metadata, and the provenance of its local proof.  The
schema does not serialize Lean proof terms.  It is intended for generators and
audits: the generated Lean source must still be elaborated and kernel-checked.

\subsection{Theorem ledger}

\path{metadata/THEOREM_LEDGER.md} classifies every main result as one of:

\begin{itemize}
\item \textbf{upstream checked}: a theorem already present in the pinned
repository;
\item \textbf{human complete}: a proof is given in this report;
\item \textbf{executable evidence}: exact reference code and tests exist;
\item \textbf{Lean candidate}: code is supplied but not compiled here;
\item \textbf{blueprint}: definitions and theorem signatures give a clear path,
with proof holes explicitly counted;
\item \textbf{conditional/application-specific}: the generic finalizer is
proved but a deep producer is not supplied.
\end{itemize}

This vocabulary is intentionally conservative.  ``Formalized'' is reserved for
kernel-checked source in the target calculus.

\subsection{Source provenance}

The source audit records:

\begin{itemize}
\item the pinned repository commit;
\item paths and blob hashes of the key upstream files;
\item the status statement in the parent iPRA README;
\item the earlier foundation packets compared;
\item the newly added direct derivation library;
\item the exact absence of Lean executables in the artifact environment.
\end{itemize}

The build does not download dependencies.  It compiles the PDF, runs the Python
tests, produces the demo JSON, and refreshes checksums.  Lean checking must be
performed in the pinned ProveIt workspace.

\subsection{PDF validation}

The report is compiled twice with \path{pdflatex}, inspected with
\path{pdfinfo}, preflighted, rendered to PNG at high resolution, and sampled
across the document for visual defects.  The validation directory records page
count, metadata, preflight output, and the pages inspected.  Intermediate
renders are not included in the final ZIP.

\subsection{Negative claims}

This package does not claim:

\begin{itemize}
\item that its new Lean files compile;
\item that the Python program is an iPRA proof checker;
\item that every proof packet in the repository is correct;
\item a universal primitive-recursive evaluator;
\item a general effective proof of Faltings, Siegel, or every stated
prime-distribution theorem;
\item a canonical selector for singular roots, spectra, quotient
representatives, or compatible infinite paths;
\item that classical finiteness by itself yields a search cutoff;
\item that a semantic Lean theorem is interchangeable with
\(\Prf\;\PRAAx\).
\end{itemize}

The positive claim is narrower and more actionable: current checked
\path{ConciseIPRA} already exposes a proof language in which a genuine generic
finite graph compiler can be very small, and the current checked direct
arithmetic library can exchange theorems with that compiler.
